\begin{corollary}[S-38: Split residual algebra for a mixed action]
\label{cor:semireg-mixed-reduced-rank}
Let \(G\subset X\times\widehat X\) be an effective finite subgroup, equipped with a coherent action \(\Phi_g=t_{a_g}^*(-)\otimes P_g\). Let a positive-virtual-rank perfect object \(\mathcal E\) have a coherent \(G\)-equivariant lift. Define its pure tensor subgroup \[
 T=\ker(G\longrightarrow X),\qquad H=G/T.
 \tag{S38.16}
\] Apply S-38(a) to the restricted \(T\)-action, giving the center isogeny \(p:B\to X\), its degree \(r\), the Azumaya degree \(d\), the reduced rank \(m=N/(rd)\), and the Brauer class \(\beta\in\operatorname{Br}(B)\).

Assume \(\beta=0\). Then the ordinary object \(F\) constructed from cohomology sheaves in S-38(b) can be given a coherent action of \(H\) by translations and tensoring with line bundles in \(\operatorname{Pic}^0(B)\). This action is effective; its translation projection is itself injective. Its normalized character is precisely (S38.9).

In the balanced real-quadratic fourfold setting, if the original \(G\)-equivariant lift is simple, gluable and fully semiregular, then \[
 k=\frac{4DN^2q^2}{|G|}\in\{2,4\}
\] and the reduced rank satisfies \[
 \boxed{\Delta_0\mid m,\qquad k=2\Longrightarrow2\mid m.}
 \tag{S38.17}
\] Consequently \textbf{reduced rank \(m=1\) is impossible for every effective finite coherent translation--tensor action}, because \(m=1\) already forces \(\beta=0\) in S-38(a). A nontrivial translation quotient cannot repair the excluded reduced-rank-one tensor case.

The argument does not claim that \(F\) is simple, gluable or semiregular. Only its actual equivariance, its normalized character, and the original numerical value \(k\) are used.
\end{corollary}
\begin{proof}
\par

\medskip\noindent\textbf{1. Taking equivariance in two stages}\quad

Write the specified coherence maps as \[
 \mu_{g,h}:\Phi_g\Phi_h\longrightarrow\Phi_{g+h}.
\] Work first in the abelian category of coherent sheaves, where all these functors are exact. For a \(T\)-equivariant sheaf \((M,\lambda_t)\) define \(A_g(M)=\Phi_gM\) and give it the \(T\)-equivariance maps \[
 \lambda_t^{(g)}
  =\Phi_g(\lambda_t)\circ\mu_{g,t}^{-1}\circ\mu_{t,g}.
 \tag{S38.18}
\] The equality \(g+t=t+g\) is used here. The pentagon identity for \(\mu\) verifies that (S38.18) is a \(T\)-equivariant structure, and that the original \(\mu_{g,h}\) induce coherent isomorphisms \(A_gA_h\simeq A_{g+h}\) on this category.

For \(t\in T\), the map \(\lambda_t\) is a natural isomorphism \[
 \epsilon_t:A_t\longrightarrow\operatorname{Id}
 \tag{S38.19}
\] on the category of \(T\)-equivariant sheaves. Its being a morphism of equivariant objects follows from \[
 \lambda_t\lambda_u^{(t)}
    =\lambda_{t+u}\mu_{u,t}
    =\lambda_u\Phi_u(\lambda_t).
\] The maps \(\epsilon_t\) satisfy the original multiplication coherence, by the defining identity for the \(\lambda_t\).

Choose representatives \(s(h)\in G\) for \(h\in H\), with \(s(0)=0\), and use \(A_{s(h)}\) as the quotient functors. In a composition, the difference \[
 t(h,k)=s(h)+s(k)-s(h+k)\in T
\] is removed by (S38.19). This defines a coherent \(H\)-action: the pentagon reduces to the pentagon for \(\mu\), the multiplication coherence for \(\epsilon_t\), and the identity \[
 t(h,k)+t(h+k,l)=t(k,l)+t(h,k+l).
\] Restricting a \(G\)-equivariant sheaf to \(T\) and using its maps for \(s(h)\) makes it an object for this quotient action. Conversely these quotient maps together with the \(T\)-equivariance recover its \(G\)-equivariance. This is an explicit construction at the abelian-sheaf level, with no comparison theorem for derived equivariant categories being used.

\par

\medskip\noindent\textbf{2. The quotient acts by translations on the center cover}\quad

The action \(A_g\) transports the algebra \(\mathscr A\) by the commutation maps underlying (S38.18). Thus it induces a map of its center cover over the translation \(t_{a_g}\) of \(X\). The canonical isomorphism (S38.19) implies that elements of \(T\) act trivially on the center. The induced automorphisms of \(B\) therefore give a genuine action \(\sigma_h\) of \(H\), and \[
 p\circ\sigma_h=t_{a_{s(h)}}\circ p.
 \tag{S38.20}
\] Each \(\sigma_h\) is a translation. Indeed \(b\mapsto\sigma_h(b)-b\) maps the connected variety \(B\) to the finite set \(p^{-1}(a_{s(h)})\), so it is constant. Write \(\sigma_h=t_{b_h}\). Composition gives a homomorphism \(h\mapsto b_h\).

This homomorphism is injective: if \(b_h=0\), (S38.20) gives \(a_{s(h)}=0\), so \(s(h)\in T\) and \(h=0\). Thus the center-cover action is a free finite translation action. This argument does not require that the original tensor commutator be nondegenerate.

\par

\medskip\noindent\textbf{3. Morita transport has only degree-zero tensor factors}\quad

Choose a splitting bundle \(S\) with \(\mathcal End(S)=\mathscr A_B\). Conjugate the quotient action of Step 1 by the ordinary Morita equivalence between \(\mathscr A_B\)-modules and coherent sheaves on \(B\). Each resulting functor covers \(\sigma_h\) and is locally the Morita equivalence of a matrix algebra with itself. Hence it has the form \[
 F\longmapsto Q_h\otimes\sigma_h^*F
 \tag{S38.21}
\] for a line bundle \(Q_h\) on \(B\). Explicitly, let \(M'_h\) be the \(\mathscr A_B\)-module corresponding to the transform of the original \(T\)-equivariant sheaf \(p_*S\) by \(\Phi_{s(h)}\). Then \(M'_h\) is locally free of rank \(d\), and \[
 Q_h=\mathcal Hom_{\mathscr A_B}(S,M'_h),\qquad
 M'_h\simeq S\otimes Q_h.
 \tag{S38.22}
\] The coherence in Step 1 transports to the functors (S38.21).

These \(Q_h\) lie in \(\operatorname{Pic}^0(B)\). On rational cohomology, translations and tensoring with \(P_{s(h)}\in\operatorname{Pic}^0(X)\) do not change Chern characters. Riemann--Roch for the finite étale map \(p\) therefore gives \[
 p_*\operatorname{ch}(M'_h)
 =\operatorname{ch}(\Phi_{s(h)}p_*S)
 =\operatorname{ch}(p_*S)
 =p_*\operatorname{ch}(S).
\] Since \(p_*:H^*(B,\mathbf Q)\to H^*(X,\mathbf Q)\) is an isomorphism, \(\operatorname{ch}(M'_h)=\operatorname{ch}(S)\). Taking degree two in (S38.22) gives \(d c_1(Q_h)=0\) in rational cohomology. The integral cohomology of an abelian variety is torsion-free, so \(c_1(Q_h)=0\) integrally, and \(Q_h\in\operatorname{Pic}^0(B)\).

The classes \((b_h,Q_h)\) form a finite subgroup of \(B\times\widehat B\): the coherence of (S38.21), together with translation invariance of classes in \(\operatorname{Pic}^0\), gives the group law. Its projection to \(B\) is the injective map in Step 2, so the action is effective.

Every \(\mathcal H^i(\mathcal E)\) is a \(G\)-equivariant coherent sheaf. Its Morita transform \(F_i\) consequently acquires the coherent quotient action (S38.21). Their direct sum, placed in the original cohomological degrees, is an actual perfect \(H\)-equivariant object \(F\) of virtual rank \(m\). The Chern-character proof of S-38 is unchanged, so \[
 e^{-c_1(F)/m}\operatorname{ch}(F)=m p^*v_q,
 \qquad D_B=rD.
 \tag{S38.23}
\]

\par

\medskip\noindent\textbf{4. Coprime descent and the final rank restriction}\quad

Put \(h=|H|\). Since \(|T|=rd^2\), \(|G|=hrd^2\) and \(N=rdm\), the original numerical identity becomes \[
 k=\frac{4DN^2q^2}{|G|}
   =\frac{4D_Bm^2q^2}{h}.
 \tag{S38.24}
\] Let \(H_0\subset H\) be the product of its primary subgroups for primes not dividing \(m\), and let \(H_m=H/H_0\). Apply the descent assertion of S-35 to the actual \(H\)-equivariant object \(F\) and the subgroup \(H_0\). That assertion requires neither simplicity nor semiregularity of \(F\). Together with S-30 it gives \[
 \frac{m^4D_Bq^2}{|H_0|}
   =\frac{m^2 k |H_m|}{4}
   \in\Delta_0\mathbf Z_{>0}.
 \tag{S38.25}
\] Every prime dividing \(|H_m|\) divides \(m\). Thus an odd prime dividing \(\Delta_0\) but not \(m\) would make the middle expression a unit at that prime, a contradiction. If \(m\) is odd, then \(|H_m|\) is odd; \(k=2\) makes that expression a half-integer, while \(k=4\) makes it odd and therefore incompatible with \(2\mid\Delta_0\). These are exactly (S38.17).

Finally \(m=1\) forces \(\beta=0\) by the twisted determinant in S-38(a), so the argument applies and contradicts \(\Delta_0>1\). This proves the all-mixed-action reduced-rank-one exclusion.
\end{proof}
